\section{Station Metadata Dataset (GBFS)}
\begin{comment}
    2) The from station\_information and from station\_status can be merged into a single list as the introduction refers to a "unified station record". 
    -> I would keep them separate to say who is from who

    3) Why operational flags (is\_installed, is\_renting, is\_returning) use a different style than the other fields?
    The style should be consistent along all the report, either all in \texttt{code} or all in \textbf{bold}. (I think the first option is better)
    -> Since referring to the actual field names, I would keep them in \texttt{code} style. Meanwhile Operational Flags is a list entry, so I would keep it in \textbf{bold} style. 

    5) The last sentences about the filtering of inactive stations and the treatment of remaining active stations are not clear to me. But we can simply remove them, 
    as they again are not relevant to the technical report.
 \end{comment}
To enrich the static station data provided by Citi Bike, we incorporate their real-time feed, available on the \href{https://www.citibikenyc.com/system-data}{\underline{Citi Bike System Data}} portal.
This live feed offers direct insight into the current capacity, operational status, and physical infrastructure of individual stations.

\subsection{Dataset Structure}
Unlike historical trip records, station metadata is served via a public JSON API adhering to MobilityData's \href{https://gbfs.org/}{\underline{GBFS}} standard.
From this service, we query two key Citi Bike endpoints:
\begin{itemize}
    \item \texttt{station\_information.json}: Static properties like location and capacity, whose content changes only when the infrastructure is modified.
    \item \texttt{station\_status.json}: Live availability and operational state.
\end{itemize}

\subsubsection{Original Structure}
The two endpoints share the \texttt{station\_id} field, on which they are merged to produce a unified station record comprising:

\begin{itemize}
    \item \textbf{From \texttt{station\_information}}:
    \begin{itemize}
        \item \textbf{Station ID}: Internal GBFS identifier, a UUID for most stations and a numeric string for the rest.
        \item \textbf{Short Name}: Public station identifier used for joins with Citi Bike trip data.
        \item \textbf{Name}: Station name.
        \item \textbf{Latitude} and \textbf{Longitude}: Station coordinates.
        \item \textbf{Capacity}: Total number of docks installed at the station.
    \end{itemize}

    \item \textbf{From \texttt{station\_status}}:
    \begin{itemize}
        \item \textbf{Bikes Available}: Currently rentable bikes.
        \item \textbf{Vehicle Types Available}: Per-type counts of the bikes currently at the station.
        \item \textbf{Docks Available}: Currently available docks.
        \item \textbf{Operational Flags}: \texttt{is\_installed}, \texttt{is\_renting}, and \texttt{is\_returning}.
    \end{itemize}
\end{itemize}

\subsubsection{Extracted Features}
Using the three operational flags (\texttt{is\_installed}, \texttt{is\_renting}, and \texttt{is\_returning}), we derive a boolean \texttt{active} indicator to filter out unavailable stations from live views.
Evaluating all three status flags is necessary because physical installation does not imply service: stations can remain installed while suspending both rentals and returns.
Relying solely on \texttt{is\_installed} would therefore misclassify inactive stations as available.

\subsection{Analysis and Insights}
This section examines the structural and operational characteristics of the Citi Bike network, providing a foundational snapshot of its system infrastructure.
Because the feed is a dynamic snapshot, this analysis represents a single point in time rather than a strictly reproducible state.

\paragraph{Capacity Distribution}
Figure~\ref{fig:station_capacity} highlights the contrast between smaller local stations and major transit hubs.
Notably, the count at $0$ docks does not indicate corrupted data; rather, it represents $42$ inactive stations currently present in the feed (as detailed in our \href{https://github.com/davide-dona/nyc-bike-data-visualization/blob/main/src/notebooks/analysis.ipynb}{\underline{Jupyter notebook}}).

\begin{figure}[ptb]
    \centering
    \includegraphics[width=\textwidth]{../media/technical/station_capacity.png}
    \caption{Distribution of station capacity, measured as the number of docks per station.}
    \label{fig:station_capacity}
\end{figure}

\paragraph{Live Availability} %#TODO: FIX
Figure~\ref{fig:station_availability} provides a comprehensive breakdown of the real-time operational state of the system.
Delivering these status classifications dynamically at the station level is precisely what makes this live feed a valuable component of the interactive dashboard.

\begin{figure}[ptb]
    \centering
    \includegraphics[width=\textwidth]{../media/technical/station_availability.png}
    \caption{System-wide capacity breakdown across all, active, and inactive stations. The left panel gives the size of each group, in stations and in docks. In the right panel, total capacity is fully decomposed into six mutually exclusive states: rentable classic bikes, rentable e-bikes, free docks, out-of-service bikes, out-of-service docks, and unreported capacity.}
    \label{fig:station_availability}
\end{figure}

\subsection{Data Quality}
To maintain data integrity and prevent visual distortion, our initial validation focused on three core quality issues: unreported capacity, feed self-consistency, and identifier alignment.
Ultimately, as confirmed in our \href{https://github.com/davide-dona/nyc-bike-data-visualization/blob/main/src/notebooks/analysis.ipynb}{\underline{analysis notebook}}, the live feed contains no missing station values.

\paragraph{Inactive Stations}
Stations undergoing installation, repair, or decommission are grouped under the \texttt{inactive} category, representing $99$ out of $2{,}462$ stations ($4.0\%$) and $1{,}950$ docks ($2.6\%$ of total capacity).
As illustrated in Figure~\ref{fig:station_availability}, these stations declare capacity while broadcasting almost no live counts, and the few bikes assigned to them are nearly all out of service.
Retaining them would distort overall capacity metrics and render ghost (false empty) stations on the interactive map without contributing usable availability.

\paragraph{Unreported Capacity}
Unreported capacity (Figure~\ref{fig:station_availability}) represents the discrepancy when a station's static capacity exceeds the sum of its live bike and dock counts. 
This arises from inconsistencies between the two endpoints, which are not guaranteed to be synchronized in the updates.
Filtering out inactive stations resolves the vast majority of these mismatches without altering real-time availability. 
For any remaining discrepancies at active stations, live bike and dock counts are prioritized, treating static capacity as a baseline.

\paragraph{Identifier Consistency}
GBFS identifies stations in two ways: an internal \texttt{station\_id} (either a UUID or a numeric string) and a public \texttt{short\_name}.
Trip history files use only the latter: $93.8\%$ of station references in historical trips match a \texttt{short\_name}, whereas none match a \texttt{station\_id}.
Therefore, all joins between live metadata and trip history rely on the \texttt{short\_name} field.

\paragraph{Live-Only Metadata}
Citi Bike provides station metadata only as a live feed, with no historical archives.
Consequently, removed stations retain their past trip records but lose all feed attributes, while newly added stations lack historical data, making them appear prematurely inactive.
This bias only affects feed-specific attributes like capacity and operational status, as each trip record independently contains its own station name and coordinates.
